% titel: Blatt 3, A3: Eduktmengen für eine Zielmenge S2Cl2
% skript: Mortimer Kap. 4.5 (Stöchiometrie, Stoffmenge)

% Teil I: Werkzeuge, kein Lösungsweg.

Soll von einem Stoff eine bestimmte Masse \emph{hergestellt} werden, rechnet
man von der Produktmenge zurück auf die dazu nötigen Eduktmengen --- über die
Stoffmenge $n$ und die Stöchiometriefaktoren einer ausgeglichenen Gleichung.

\begin{werkzeugkasten}[Zurückrechnen von Produkt auf Edukt]
\begin{enumerate}
  \item Reaktionsgleichung ausgleichen; für jedes Edukt den Faktor $\nu$ ablesen.
  \item Stoffmenge des Produkts aus Zielmasse und molarer Masse:
        $n_{\mathrm{Prod}} = m/M$.
  \item Jedes Edukt über die Stöchiometrie: $n_{\mathrm{Edukt}} =
        \nu_{\mathrm{Edukt}}\cdot n_{\mathrm{Prod}}$.
  \item Masse jedes Edukts mit $m = n\cdot M$.
\end{enumerate}
Beispiel: $a\,\mathrm{A} + b\,\mathrm{B} \rightarrow c\,\mathrm{C}$.
Für $n_{\mathrm{C}}$ Ziel braucht man $n_{\mathrm{A}} = (a/c)\,n_{\mathrm{C}}$
und $n_{\mathrm{B}} = (b/c)\,n_{\mathrm{C}}$. Die Gesamtmasse der Edukte ist
so bemessen, dass sie der Zielmasse des Produkts (plus nichts Verlorenes)
entspricht.

\textit{Skript:} Mortimer Kap. 4 (Stöchiometrie, Molmasse), Übung 3/3.
\end{werkzeugkasten}

\begin{achtung}[Darauf achten]
\begin{itemize}
  \item Die Schreibweise von $\mathrm{S_2Cl_2}$ sauber lesen: zwei Schwefel-
        und zwei Chlor-Atome pro Molekül; die Stöchiometriefaktoren daher
        $2$ bzw.\ $1$ (für $\mathrm{Cl_2}$).
  \item Produkt zuerst in $n$ umrechnen; nicht Masse $\times$ Faktor direkt,
        weil sich die molaren Massen unterscheiden.
  \item Für Chlor den Faktor auf $\mathrm{Cl_2}$ beziehen ($1\ \mathrm{Cl_2}$
        pro $\mathrm{S_2Cl_2}$), nicht auf einzelne Cl-Atome.
  \item Einheiten g/mol sauber notieren; Zwischenwerte nicht abrunden.
\end{itemize}
\end{achtung}
